\diarytitle{0032}{定数係数2階線形非同次微分方程式の一般解（非同次項が多項式の場合）}

未定係数法を用いて特殊解を求める。定数係数2階線形非同次方程式 $y'' + a y' + b y = f(x)$ の一般解は、対応する同次方程式の一般解 $y_h$ と、非同次方程式の特殊解 $y_p$ の和
\[
y = y_h + y_p
\]
で与えられる。特殊解は $f(x)$ の形に応じて多項式・指数関数・三角関数などの仮定形を置き、方程式に代入して未定係数を決定する。ただし仮定した形が同次解と重なる場合（共鳴）は、仮定形に $x$ の適当なべきを掛けて同次解と一次独立にする必要がある。

\medskip
\noindent
\textbf{同次解:} 特性方程式 $r^{2} - 1 = 0$ より $r = \pm 1$。よって
\[
y_h = C_{1}e^{x} + C_{2}e^{-x}
\]

\noindent
\textbf{特殊解:} $f(x) = x^{2}+1$ は2次多項式で、同次解に多項式項は含まれないので共鳴なし。$y_p = Ax^{2}+Bx+C$ とおくと $y_p'' = 2A$ より
\[
y_p'' - y_p = 2A - (Ax^{2}+Bx+C) = -Ax^{2} - Bx + (2A - C)
\]
これが $x^{2}+1$ に等しいので
\[
-A = 1,\qquad -B = 0,\qquad 2A - C = 1
\]
よって $A=-1,\ B=0,\ C = 2A-1 = -3$。すなわち $y_p = -x^{2}-3$。

\noindent
したがって一般解は
\[
\boxed{\; y = C_{1}e^{x} + C_{2}e^{-x} - x^{2} - 3 \;}
\qquad (C_1, C_2 \text{ は任意定数})
\]

（検算: $y_p''=-2$ より $y_p''-y_p = -2-(-x^2-3) = x^2+1$ OK）
